Wer sind die Adressaten der Rechnungslegung?

Interne Adressaten (Management)
Externe Adressaten 
Gläubiger, Aktionäre, Arbeitnehmer, Finanzbehörden, Lieferanten, Kunden, Aufsichtsbehörden

Interessen dieser Adressaten sind unter anderem
Management: Planungsdaten, Erfolgsmessung, Einsparpotenziale 
Gläubiger: Wie sicher sind Kredite?
Aktionäre: Wie hoch sind Gewinne/Verluste
Arbeitnehmer: Wie sicher ist der Arbeitsplatz/das Gehalt?
Finanzbehörden: Wie hoch sind die zu zahlenden Steuern?
Lieferanten: Wird die Ware bezahlt?
Kunden: Fällt die Warenlieferung kurzfristig wegen Insolvenz aus, sind Garantieansprüche
durchsetzbar
Aufsichtsbehörden: Sicherung der Interessen der Verbraucher (Versicherungsnehmer)



Was sind die Bestandteile des externen Rechnungswesens
Finanzbuchhaltung
Jahresabschluss (Bilanz und GuV erweitert um Anhang)
Lagebericht

Zu welchem Detailgrad dies geschieht ist abhängig von der Organisationsform, 
(reregelt nach (§242 Abs. 3 HGB), (§264 Abs. 1 HGB), §297 HGB) hier aufsteigend gelistet
(spätere müssen alles vom Vorgänger erfüllen):
Kaufmann
-Bilanz
-Gewinn und Verlustrechnung
Kapitalgesellschaft
- Anhang
- Lagebericht
Kapitalmarktorientierte Kapitalgesellschaft und Konzerne
- Kapitalflussrechnung
- Eigenkapitalspiegel
- Optional Segmentberichterstattung

Was ist die Handelsrechtliche Buchführungspflicht gem. §238 Abs. 1 HGB

Jeder Kaufmann ist verpflichtet, Bücher zu führen und
in diesen seine Handelsgeschäfte und die Lage seines
Vermögens nach den Grundsätzen ordnungsmäßiger
Buchführung ersichtlich zu machen.
Die Buchführung muss so beschaffen sein, dass sie
einem sachverständigen Dritten innerhalb
angemessener Zeit einen Überblick über die Geschäftsvorfälle und über die Lage des Unternehmens
vermitteln kann.
Die Geschäftsvorfälle müssen sich in ihrer Entstehung
und Abwicklung verfolgen lassen.



Was ist die Bedeutung vom Anhang im Jahresabschluss (keine Klausur bisher)

Gesetzliche Grundlage:
Die gesetzlichen Vertreter einer Kapitalgesellschaft haben den Jahresabschluss (§242) um einen Anhang zu erweitern, der mit der Bilanz und der Gewinn- und Verlustrechnung eine Einheit bildet… (§264 Abs. 1 HGB)
Informationsfunktion 
Der Anhang hat die Aufgabe, zusammen mit Bilanz sowie Gewinn- und Verlustrechnung unter Beachtung der Grundsätze ordnungsmäßiger Buchführung ein den
tatsächlichen Verhältnissen entsprechendes Bild der Vermögens-, Finanz- und Ertragslage zu vermitteln.

Entlastungsfunktion:
Zwecks übersichtlicher Gestaltung von Bilanz und Gewinn- und Verlustrechnung können Angaben auch im Anhang gegeben werden. 
Die Gleichstellung des Anhangs mit der Bilanz und der Gewinn- und Verlustrechnung im Rahmen des Jahresabschlusses erlaubt es, 
ohne Informationsverlust Angaben in den Anhang zu übernehmen,
die sonst in einem der beiden anderen Rechnungslegungsinstrumente zu machen wären; diese Teile des Jahresabschlusses werden dadurch entlastet.

Erläuterungsfunktion: 
Hierzu zählen alle diejenigen Informationen des Anhangs, welche Bilanz und Gewinn- und Verlustrechnungsposten kommentieren oder interpretieren. Neben
Einzelangaben zu den in Bilanz und Gewinn- und Verlustrechnung offenlegungspflichtigen Posten zählen hierzu u. a. die Angabe von Bilanzierungs- und
Bewertungsmethoden (§284 Abs. 2 Nr. 1 HGB), oder die Angabe und Begründung von Methodenänderungen (§284 Abs. 2 Nr. 2 HGB).

Ergänzungsfunktion: Die Erläuterungsfunktion wird insbesondere durch die Ergänzungsfunktion verstärkt, die sich –im Gegensatz zur Erläuterungsfunktion– auf Informationen erstreckt, die
weder in Bilanz noch in Gewinn- und Verlustrechnung (bzw. bei nicht zur Aufstellung eines Konzernabschlusses verpflichteten kapitalmarktorientierten
Kapitalgesellschaften in Bilanz, Gewinn- und Verlustrechnung, Kapitalflussrechnung und Eigenkapitalspiegel) gegeben werden. Hierzu zählen insbesondere die Angabepflichten nach §285 Nr. 7, 10, 11 a, 15, 15a, 16, 19, 21, 26, 31, 32, 33 und 34 HGB.

Korrekturfunktion: 
Sofern Bilanz und Gewinn- und Verlustrechnung in besonderen Fällen nicht in der Lage sind, ein den tatsächlichen Verhältnissen entsprechendes 
Bild der Vermögens-, Finanz- und Ertragslage zu vermitteln, sind im Anhang gemäß §264 Abs. 2 Satz 2 HGB zusätzliche Angaben zu machen, 
um in der Gesamtheit des Jahresabschlusses
die Anforderung der Generalnorm des §264 Abs. 2 Satz 1 HGB zu erfüllen.

Spezielle Angaben bei Versicherungen:
• Angaben zu versicherungstechnischen Posten, z.B. Art und Höhe des Abwicklungsergebnisses der Schadenrückstellung, falls dieses erheblich ist
(§41 Abs. 5 RechVersV)
• Angaben zu Kapitalanlagen, z.B. Zeitwerte der Kapitalanlagen (§54 RechVersV)

Was ist die Bedeutung vom Lagebericht im Jahresabschluss (keine Klausur bisher)

Gesetzliche Grundlage: 
Die gesetzlichen Vertreter einer Kapitalgesellschaft haben den Jahresabschluss (§242 HGB) um einen Anhang zu erweitern, 
der mit der Bilanz und der Gewinn- und Verlustrechnung eine Einheit bildet, sowie einen Lagebericht aufzustellen (§264 Abs. 1 HGB).
Im Lagebericht sind der Geschäftsverlauf einschließlich des Geschäftsergebnisses und die Lage der Kapitalgesellschaft so darzustellen, dass ein 
den tatsächlichen Verhältnissen entsprechendes Bild vermittelt wird. Er hat eine ausgewogene und umfassende, dem Umfang und der Komplexität der Geschäftstätigkeit entsprechende Analyse des Geschäftsverlaufs und der Lage der Gesellschaft zu enthalten (§289 Abs. 1 HGB).

Geschäftsverlauf: 
Die Darstellung greift die Zahlen aus der Bilanz und GuV auf und erläutert die Ursachen. Dies beinhaltet einen Überblick über den Konzern, die Ziele der
Unternehmensleitung und über den Geschäftsverlauf sowie Veränderungen in der Ertrags- und Aufwandsstruktur, Finanzmanagement, Kapitalstruktur und Liquidität (auch außerbilanzielle Finanzierungsinstrumente).

Chancen: 
Die Chancen leiten sich aus der Unternehmensstrategie und dem Marktumfeld ab (z.B. Spartenwachstum, Kostensenkung, Ausbau Marktanteil etc.).

Risiken: 
Die Risiken leiten sich aus dem Risikomanagement des Unternehmens ab.
(z.B. versicherungstechnische Risiken, Zinsänderungsrisiken, Kreditrisiken, Liquiditätsrisiko etc.)

Prognose: 
Die Prognose leitet sich aus dem Geschäftsverlauf, den Risiken und Chancen ab. Für den Geschäftsverlauf und die Prognose sind jeweils die wesentlichen
Steuerungsgrößen zu nennen (Key Performance Indicators(KPI)). Für jeden KPI sind in der Prognose Richtung und Intensität zu nennen
(z.B. steigt stark oder sinkt moderat).

Übernahmerelevante Angaben:
Zusammensetzung des gezeichneten Kapitals, Stimmrechtsbeschränkungen, Sonderrechte, Stimmrechtskontrollen, Ausgabe-/Rückkaufbefugnisse des Vorstands,
Entschädigungsvereinbarungen bei Übernahmeangeboten (§289 a HGB) (nur für kapitalmarktorientierte Unternehmen).

Nichtfinanzielle Erklärung: 
Erläuterungen zu Umweltbelangen, Arbeitnehmerbelangen, Sozialbelangen, Achtung der Menschenrechte, Bekämpfung von Korruption (§289 b und c HGB) (nur für
kapitalmarktorientierter Unternehmen, Versicherungen und Kreditinstitute mit jeweils mehr als 500 Arbeitnehmern).

Erklärung zur Unternehmensführung:
Entsprechungserklärung zum Corporate Governance Kodex nach §161 AktG, relevante Angaben zu Unternehmensführungspraktiken, Beschreibung der Arbeitsweise
von Vorstand und Aufsichtsrat, Zusammensetzung von Vorstand und Aufsichtsrat, Erläuterung eines Diversitätskonzepts (nur für kapitalmarktorientierte Unternehmen).



Nennen sie eine Auswahl von allgemeinen Bewertungsgrundsätzen

Das wichtigste ist das nach §252 Abs. 1 Nr. 4 HGB
Vorsichtsprinzip:
    Es ist vorsichtig zu bewerten, namentlich sind alle vorhersehbaren Risiken und Verluste, die bis zum Abschlussstichtag
    entstanden sind, zu berücksichtigen, selbst wenn diese erst zwischen dem Abschlussstichtag und dem Tag der Aufstellung
    des Jahresabschlusses bekanntgeworden sind; Gewinne sind nur zu berücksichtigen, wenn sie am Abschlussstichtag
    realisiert sind.“

Daraus ergibt sich
- Vermögensgegenstände
   Von zwei möglichen Werten ist der niedrigere anzusetzen
   (Niederstwertprinzip)
   - Anlagevermögen
     (gemildertes Niederstwertprinzip)
   - Umlaufvermögen
     (strenges Niederstwertprinzip)
- Rückstellungen und Verbindlichkeiten
   Von zwei möglichen Werten ist der höhere
   anzusetzen (Höchstwertprinzip) 

Die Varianten des Niederstwertprinzip sind im Detail  ist weiterhin zu beachten
strenges Niederstwertprinzip:
- Dauerhafte Wertminderung
- Vorübergehende Wertminderung
in beiden Fällen gilt die Abschreibungspflicht 
Gemildertes Niederstwertprinzip
- Dauerhafte Wertminderung -> Abschreibungspflicht 
- Vorübergehende Wertminderung -> Abschreibungswahlrecht wenn es eine Finanzanlage ist

In Ausnahmefällen findet das Wertaufholungsgebot nach §253 Abs. 5 HGB Anwendung, z.B.
Wenn Grund für niedrigeren Wertansatz nicht mehr bestehen.

Desweiteren gibt es die "soften"
Informationsprinzipien-GoB
• Adressatenorientierung
• Entscheidungsrelevanz
• Klarheit
• Darstellungsstetigkeit
• Wesentlichkeit


Was sind die Begriffe zur Kategorisierung von Sachverhalten in der Bilanz (Strömungsgrößen)?

Aufwand/Ertrag: 
Information für externe Bilanzadressaten
• Diese beiden Begriffe sind von zentraler Bedeutung für die Buchhaltung und die Ermittlung des Jahresüberschusses
• Erträge sind der Wert der erbrachten Leistungen (z.B. Warenverkauf, Umsatzerlöse)
• Aufwendungen sind der Wert der dafür verbrauchten Güter/Dienstleistungen (z.B. Personalkosten, Verbrauch von Rohstoffen zur Leistungserstellung)
• In der Gewinn- und Verlustrechnungen werden nach Aufwand und Ertrag gegenübergestellt, um den Jahresüberschuss zu ermitteln

Einzahlung/Auszahlungen
• Information für die Finanzplanung und Investitionsplanung
• Beschreibt den Zufluss/Abfluss liquider Mittel
• Z.B. Banküberweisung an einen Lieferanten i.H.v. EUR 1.000 (Auszahlung)
• Z.B. Eingang Banküberweisung durch einen Kunden i.H.v. EUR 1.000 (Einzahlung)

Einnahme/Ausgabe
• Das Geldvermögen erhöht sich durch Einnahmen und verringert sich durch Ausgaben
• Das Geldvermögen ist definiert als Forderungen, Bankguthaben, Kasse abzgl. Verbindlichkeiten
• Die Leistungserbringung an einen Kunden führt mit Versand der Rechnung (Buchung einer Forderung) zu Einnahmen. Der Eingang einer Rechnung eines
Lieferanten und die Buchung der entsprechenden Verbindlichkeit führt zu einer Ausgabe.

Kosten/Erlöse
Informationen für die interne Unternehmenssteuerung (Controlling)
• Damit lässt sich betriebsbezogen feststellen, wie teuer ein Produkt in der Herstellung ist (Kosten) und wie hoch der Umsatz mit diesem Produkt ist (Erlös).
• Die Abgrenzung zu Aufwand/Ertrag ist, dass Kosten/Erlöse nur betriebsbedingtes Vermögen, also das Betriebsergebnis umfassen und Aufwand/Ertrag das
Gesamtergebnis. So enthalten Kosten und Erlöse auch kalkulatorisch Kosten (z.B. fiktive Abnutzung des Gebäudes oder der Maschinen). Die Aufwendungen
hingegen enthalten Abschreibungen, die nicht immer der tatsächlichen Abnutzung, sondern einem linearen Verlauf entsprechen).

Geben Sie Beispiele von Sachverhalten für die jeweiligen Strömungsgrößen

Beispiele:
Auszahlung die keine Ausgabe ist:
Bezahlung einer bestehenden Verbindlichkeit aus der Vorperiode
Auszahlung = Ausgabe:
Einkauf von Rohstoffen, Barzahlung
Ausgabe die keine Auszahlung ist:
Zieleinkauf von Rohstoffen (jetzt kaufen später zahlen)
Ausgabe die kein Aufwand ist:
Kauf von Rohstoffen, kein Verbrauch in der Periode (Lagerung)
Ausgabe = Aufwand:
Kauf und Verbrauch von Rohstoffen in einer Periode
Aufwand der keine Ausgabe ist:
Verbrauch von Rohstoffen ab Lager (Kauf in Vorperiode)
Neutraler Aufwand: 
Spenden, Verlust aus Wertpapierverkauf
Aufwand = Kosten: (Zweckaufwand/Grundkosten)
Rohstoffeinsatz
Zusatzkosten:
alkulatorische Miete


Einzahlung die keine Einnahme ist:
Kunde bezahlt Ausgangsrechnung aus der Vorperiode
Einzahlung = Einnahme:
Barverkauf von Erzeugnissen
Einnahme die keine Einzahlung ist:
Zielverkauf von Erzeugnissen (jetzt verkaufen später Entgelt erhalten)
Einnahme die kein Ertrag ist:
Kunde leistet Anzahlung für Erzeugnisse, die in der folgenden Periode erstellt werden
Einnahme = Ertrag:
Verkauf von in der laufenden Periode erstellten Erzeugnissen
Ertrag der keine Einnahme ist:
Herstellung von Erzeugnissen und Lagerung
Neutraler Ertrag:
Gewinn aus Wertpapierverkauf
Ertrag = Leistung (Zweckertrag/Grundleistung):
Verkauf von Erzeugnissen
Zusatzleistung:
unentgeltlich abgegebene Leistungen oder Selbsterstelltes Know-How


Nennen Sie grundlegende Objekte in der Buchführung 
Doppelte Buchführung:
- Verpflichtet das Unternehmen zur Kontierung seiner Geschäftsvorfälle auf mindestens zwei Konten: Konto und Gegenkonto
Geschäftsvorfall:
- Tätigkeit bzw. Vorgang, der zur Änderung der Höhe und/oder Struktur des Vermögens, der Schulden oder des Reinvermögens beiträgt
- Jede Buchung eines Geschäftsvorfalls basiert auf einem Beleg (z.B. Rechnung)
Grundbuch:
- Im Grundbuch werden alle Geschäftsvorfälle in zeitlicher Reihenfolge festgehalten
Hauptbuch:
- Im Hauptbuch werden die Geschäftsvorfälle nach sachlichen Kriterien gegliedert (Anlagevermögen, Umlaufvermögen, Eigenkapital, Fremdkapital).
- Das Hauptbuch enthält alle Sachkonten, aus denen sich der Jahresabschluss ergibt
Konto:
- Jeder Geschäftsvorfall wird auf Sachkonten gebucht (Bestands-/Erfolgskonten)
- Die Konten werden nach dem Prinzip Soll an Haben bebucht


Stellen sie Beziehungen von den "fundamentalen" Konten dar:

Konten werden mit einem
Eröffnungsbilanzwert
(Bilanzwert des Vorjahres, in der GuV
Null) eröffnet
• Dann werden alle Geschäftsvorfälle des
Jahres gebucht. Diese erhöhen oder
vermindern den Kontenbestand
• Am Jahresende wird das Konto
abgeschlossen und der verbleibende
Saldo in die Schlussbilanz/-GuV
gebucht
• Die GuV-Konten haben keinen
Eröffnungsbilanzwert
(startet immer bei Null)


Erklären sie den Enstehungsprozess der Bilanz

1. Man fasst die einzelnen Positionen des Inventars zu größeren Gruppen zusammen.
   Dabei fallen die genauen Spezifizierungen weg. Ist also zum Beispiel im Inventar
   jede Maschine nach Typ und Baujahr einzeln aufgeführt, so wird nun eine Position
   „Anlagen“ oder „Maschinen“ gebildet
2. Eine solche Zusammenfassung bedingt bereits den zweiten Schritt: den Wegfall der Mengenangaben.
3. Man ordnet Vermögen und Schulden nicht mehr hintereinander an, sondern
   stellt sie einander gegenüber, wobei traditionell die Vermögenswerte links und die
   Schuldenwerte rechts notiert werden.
4. Schließlich bildet man die bereits oben genannte Differenz, den "Saldo", zwischen
   Vermögen und Schulden, das sogenannte Eigenkapital. Dieses setzt man auf
   die kleinere Seite der Rechnung, also in der Regel auf die Seite der Schulden (des
   Fremdkapitals). Meist steht das Eigenkapital auf der rechten Seite an erster Stelle,
   obwohl es als letzter Posten durch Saldierung (Differenzbildung) von Vermögen
   und Schulden (Fremdkapital) ermittelt wird. Ist das Eigenkapital negativ so ist das Unternehmen insolvent.

Grundsätzliche Struktur ist dass auf der einen Seite die konkreten Vermögensgüter aufgezeichnet werden 
Auf der anderen Bilanzseite stehen aber nicht andere Güter, sondern der Wert derselben Vermögensgüter, nun aber unter dem Gesichtspunkt 
der Eigentumsverhältnisse betrachtet. Diese Posten fasst man unter der Bezeichnung Kapital zusammen.

Da man das Vermögen auch als „Aktiva“, das Kapital als „Passiva“ bezeichnet, gilt
ebenfalls:
\( \text{Aktiva} = \text{Passiva} \).
Bei Zerlegung des gesamten Kapitals in Eigenkapital (EK) und Fremdkapital (FK)
ergibt sich:
\( V = \text{EK} + \text{FK} \)

Was ist die Gliederung der Aktiva und wie ist die Reihenfolge in der Bilanz?

Aktiva:
Anlagevermögen (AV)
    - Sachanlagen und immaterielle Anlagenwerte
    -Finanzanlagen
Umlaufvermögen (UV)

Zum Anlagevermögen gehören alle Vermögensarten, die am Bilanzstichtag dazu
bestimmt sind, längere Zeit in der Unternehmung zu verbleiben.
Das Anlagevermögen eine Nutzung durch Gebrauch erfährt, also eine mehrfache Verwen-
dung desselben möglich ist.

Als Sachanlagevermögen sind zum Beispiel bebaute und unbebaute Grundstü-
cke, Maschinen, Betriebs- und Geschäftsausstattung anzusehen.
Immaterielle Vermögensgegenstände sind Konzessionen, Patente, Lizenzen usw.

Zum Finanzanlagevermögen gehören Anteile und Beteiligungen an anderen
Unternehmungen, Wertpapiere, die nicht nur vorübergehend von der Unternehmung
gehalten werden, sowie vom Unternehmen gewährte langfristige Kredite.

Das Umlaufvermögen, das aus Rohstoffen, Waren, eigenen Erzeugnissen, For-
derungen, Guthaben bei Kreditinstituten, Kassenbestand und ähnlichen Positionen
besteht, erfasst alle Vermögensgüter, die nicht langfristig dem Geschäftsbetrieb die-
nen sollen, sondern in Umsatzakten weiterveräußert bzw. umgeschlagen zu werden
bestimmt sind.

Die linke Seite der Bilanz ist primär funktional gegliedert, d.h. die Positionen
werden aufgrund des Verwendungszwecks abgegrenzt.

Sekundär werden die Positionen auf der linken Seite der Bilanz nach der Liquidi-
tät gegliedert und angeordnet, wobei man im Anschluss an die immateriellen Vermö-
gensgegenstände mit den als am wenigsten liquide angesehenen Grundstücken und
Gebäuden beginnt und bis zu den Kassenbeständen als den liquidesten Vermögens-
teilen gelangt.

em Vermögen steht das Kapital als dessen abstrakte Entsprechung auf der rechten
Bilanzseite (Passiva) gegenüber. 
Auch das Kapital ist nach Arten gegliedert. 
Man unterscheidet zunächst nach den Anspruchseignern in Eigenkapital und Fremdkapital.

Außerdem werden bestimmte Finanzierungsarten voneinander abgegrenzt (zum
Beispiel im Fremdkapital u.a. Anleihen, Verbindlichkeiten gegenüber Kreditinstituten,
erhaltene Anzahlungen). 

Schließlich wird das Kapital nach der Fristigkeit geordnet,


Beschreiben sie die grundlegende Buchung von Geschäftsvorfällen/Sachverhalte (Buchungssätze)

Auf den als Aktiv- oder Passivkonto eröffneten Konten werden dann während des lau-
fenden Geschäftsjahres Vorfälle, die einen Bestand mehren, auf derselben Seite
gebucht wie der Anfangsbestand, und Vorfälle, die den Bestand mindern, auf der
Gegenseite.

In der Fachsprache heißt die linke Seite eines Kontos stets Sollseite, die
rechte Seite Habenseite.

Als Grundprinzip der doppelten Buchhaltung ist also festzuhalten: Jeder zu
buchende Geschäftsvorfall wird wertmäßig im Soll und in derselben Höhe
im Haben erfasst. Bei der Sollbuchung spricht man oftmals auch davon, dass ein
Konto „belastet“ wird, bei der Habenbuchung davon, dass ein Konto „erkannt“ wird.

Die Anweisung zur Buchung erfolgt im Buchungssatz. Dabei wird
immer zuerst (also links im Buchungssatz) das Konto genannt, bei dem
im Soll (also auch links) gebucht wird, an zweiter Stelle (also rechts im
Buchungssatz) das Konto, bei dem im Haben (ebenfalls rechts) gebu-
cht wird. Als drittes ist schließlich noch der zu buchende Betrag zu
nennen. Verknüpft werden die beiden Konten durch ein „an“.19
E.g.
Technische Anlagen
und Maschinen

an

Verbindlichkeiten

1.000



Malen Sie den Haben/Soll für Aktiv/Passivkonten auf

INSERT IMAGE FROM PAGE 24


Geben sie die 4 Typen von Buchungen/Bilanzwirkungen an

Aktivtausch:
Die Buchung führt zu einer Mehrung des Wertes bei einem Vermögenskonto (Soll-
buchung auf einem Aktivkonto) und gleichzeitig zu einer ebenso hohen Minde-
rung des Wertes bei einem anderen Vermögenskonto (Habenbuchung auf einem
Aktivkonto). Die Summe der Aktivseite und auch der Passivseite der Bilanz (Bilanz-
summe) bleibt unverändert.

INSERT IMAGE FROM PAGE 26

Passivtausch:
Die Buchung führt zu einer Mehrung des Wertes bei einem Kapitalkonto (Habenbu-
chung auf einem Passivkonto) und gleichzeitig zu einer ebenso hohen Minderung
des Wertes bei einem anderen Kapitalkonto (Sollbuchung auf einem Passivkonto).
Die Bilanzsumme bleibt auch hier unverändert.

INSERT IMAGE FROM PAGE 26 

Bilanzverlängerung:
Die Buchung führt zu einer Mehrung des Wertes bei einem Vermögenskonto (Soll-
buchung auf einem Aktivkonto) und gleichzeitig zu einer ebenso hohen Mehrung
des Wertes bei einem Kapitalkonto (Habenbuchung auf einem Passivkonto). Die
wertmäßige Summe der Aktivseite und die der Passivseite der Bilanz werden um
den gleichen Betrag vergrößert.20 Die Bilanz ist nach wie vor im Gleichgewicht.

INSERT IMAGE FROM PAGE 26 

Bilanzverkürzung:
Die Buchung führt zu einer Minderung des Wertes bei einem Vermögenskonto
(Habenbuchung auf einem Aktivkonto) und gleichzeitig zu einer ebenso hohen
Minderung des Wertes bei einem Kapitalkonto (Sollbuchung auf einem Passiv-
konto). Die wertmäßige Summe der Aktiv- und Passivseite der Bilanz wird gleich-
mäßig verringert, die Bilanz bleibt im Gleichgewicht.


INSERT IMAGE FROM PAGE 26 


Geben sie Beispiele für Geschäftsvorfälle und eine entsprechende Ansicht der Buchungssätze (1)

Einstellung des Jahresüberschusses i.H.v. EUR 15.000 aus dem Jahr x0 
(der Jahresüberschuss des Jahres x0 wurde in der Eröffnungsbilanz des Jahres x1 auf das Konto
Gewinnvortrag vorgetragen) in die andere Gewinnrücklage

1. Gewinnvortrag EUR 15.000 an Andere Gewinnrücklage EUR 15.000

INSERT IMAGE FROM SLIDE 40


Tilgung eines Kredites in Höhe von EUR 35.000
    2. Verbindlichkeiten EUR 35.000 an Bank EUR 35.000

INSERT IMAGE FROM SLIDE 41


Zinszahlung für Kredite in Höhe von EUR 50.000
3. Zinsaufwand EUR 50.000 an Bank EUR 50.000

INSERT IMAGE FROM SLIDE 42

Für einen neuen Rechtsstreit wird eine Rückstellung i.H.v. EUR 10.000 gebildet
5. Sonstiger betrieblicher Aufwand (sbA) EUR 10.000 an Rückstellungen EUR 10.000


Kreditaufnahme über EUR 100.000
Bank EUR 100.000 an Verbindlichkeit EUR 100.000


INSERT IMAGE FROM SLIDE 43


Für einen neuen Rechtsstreit wird eine Rückstellung i.H.v. EUR 10.000 gebildet
Sonstiger betrieblicher Aufwand (sbA) EUR 10.000 an Rückstellungen EUR 10.000

INSERT IMAGE FROM SLIDE 44




Geben sie Beispiele für Geschäftsvorfälle und eine entsprechende Ansicht der Buchungssätze (2)

Die Immateriellen Vermögensgegenstände werden mit EUR 20.000 außerplanmäßig abgeschrieben
Außerplanmäßige Abschreibung EUR 20.000 an Immaterielle Vermögensgegenstände EUR 20.000 
(die Buchung erfolgt hier vereinfacht auf dem Konto Abschreibung, es könnte auch
ein separates Unterkonto für außerplanmäßige Abschreibungen bestehen)
Auf dem Konto Immaterielle VG gibt es keine weiteren Buchungen, daher wird der Kontoabschluss hier gezeigt


INSERT IMAGE FROM SLIDE 45


Es werden folgende Sachanlagen neu angeschafft und per Bank bezahlt:
7. Auto EUR 20.000
8. Maschine EUR 60.000
9. Möbel EUR 10.000

7. Sachanlagen (KFZ) EUR 20.000 an Bank EUR 20.0008.
8. Sachanlagen (Maschine) EUR 60.000 an Bank EUR 60.0009.
9. Sachanlagen (Betriebs-und Geschäftsausstattung = BGA) EUR 10.000 an Bank EUR 10.000


Die Sachanlagen werden zum Jahresende planmäßig mit EUR 30.000 abgeschrieben
10. Abschreibung EUR 30.000 an Sachanlagen EUR 30.000

Es wird Ware an 3 Kunden auf Ziel verkauft,
Kunde 1 EUR 100.000
Kunde 2 EUR 250.000
Kunde 3 EUR 50.000

11. Forderung EUR 100.000 an Umsatzerlöse EUR 100.000
12. Forderung EUR 250.000 an Umsatzerlöse EUR 250.000
13. Forderung EUR 50.000 an Umsatzerlöse EUR 50.000
Auf dem Konto Immaterielle VG gibt es keine weiteren Buchungen, daher wird der Kontoabschluss hier gezeigt

Es erfolgen 3 Einzahlungen zu bestehenden Forderungen in Höhe von
Einzahlung 1 EUR 50.000
Einzahlung 2 EUR 250.000
Einzahlung 3 EUR 50.000

14. Bank EUR 50.000 an Forderung EUR 50.000
15. Bank EUR 250.000 an Forderung EUR 250.000
16. Bank EUR 50.000 an Forderung EUR 50.000
Auf beiden Konten gibt es keine weiteren Buchungen, daher wird der Kontoabschluss hier jeweils gezeigt


